\documentclass[11pt,a4paper]{article}

% --- Essential Packages ---
\usepackage[utf8]{utf8}
\usepackage[margin=1in]{geometry}
\usepackage{amsmath,amssymb}
\usepackage{booktabs}
\usepackage{hyperref}
\usepackage{caption}
\usepackage{listings}
\usepackage{xcolor}
\usepackage{setspace}
\usepackage{enumitem}

\onehalfspacing

% --- Hyperref Setup ---
\hypersetup{
    colorlinks=true,
    linkcolor=navyblue!80!black,
    citecolor=navyblue!80!black,
    urlcolor=blue
}
\definecolor{navyblue}{rgb}{0.0, 0.0, 0.5}

% --- XML Listing Style ---
\lstset{
    language=XML,
    basicstyle=\ttfamily\footnotesize,
    columns=fullflexible,
    breaklines=true,
    frame=single,
    backgroundcolor=\color{gray!8},
    keywordstyle=\color{navyblue},
    commentstyle=\color{gray}
}

% --- Title Information ---
\title{\Large \textbf{The Algorithmic Stage: Software Shepherding, Telemetry Control, and the Paradox of Spectatorship in the Quantified Workplace}}

\author{\textbf{Vivien Jiaqian Zhu} \\
Department of East Asian Languages and Cultures \\
University of California, Berkeley \\
\texttt{vjzhu@berkeley.edu}}

\date{\today}

\begin{document}

\maketitle

\begin{abstract}
This paper develops an integrated theoretical and empirical framework for analyzing workplace surveillance in remote and hybrid labor environments. By synthesizing Brett Aho’s concepts of \textit{software shepherding} and \textit{telemetry control} with Vivien Jiaqian Zhu’s dramaturgical model of the \textit{corporate stage} and the \textit{paradox of spectatorship}, we account for the dual nature of contemporary algorithmic management. We demonstrate that algorithmic control operates through a closed cybernetic loop comprising two distinct registers: an outer architectural circuit that enforces code-level constraints and eliminates worker workarounds, and an inner subjective circuit that compels workers to internalize those constraints through performative metric curation. Background telemetry acts as the connective tissue between these circuits, transforming routine workplace activities into real-time data streams that feed both software enforcement engines and worker-facing productivity dashboards. We outline a bimodal empirical research agenda combining critical software auditing with digital ethnography to examine how this dual-circuit architecture produces preemptive self-shepherding, and consider its expansion under generative AI workplace systems.
\end{abstract}

\noindent\textbf{Keywords:} Algorithmic Management, Software Shepherding, Corporate Stage, Telemetry Control, Paradox of Spectatorship, Remote Labor, Digital Ethnography.

\clearpage

% --- Section 1 ---
\section{Introduction and Theoretical Convergence}

The rapid normalization of remote and hybrid work environments has transformed the spatial, structural, and psychological organization of contemporary labor. Far from liberating workers from physical oversight, the migration of routine tasks into enterprise platforms (e.g., Microsoft Teams, Slack, Zoom, Teramind) has embedded labor within a sensor-saturated digital habitat. Within these environments, managerial authority is rarely exercised through direct human observation; instead, it is mediated by automated tracking systems, passive background telemetry, and real-time algorithmic evaluations. 

To conceptualize how control functions in these environments, critical scholars must move beyond traditional models of panoptic surveillance. Classic panopticism relies on the \textit{possibility} of an unseen human observer watching from a central tower, inducing discipline through uncertainty. Modern algorithmic management, however, functions through continuous, automated data ingestion and real-time interface modification. It does not merely watch labor; it actively structures, limits, and evaluates the parameters within which labor can occur.

To account for this double-sided operation, this paper synthesizes two complementary theoretical frameworks: Brett Aho’s STS-grounded theory of \textbf{software shepherding} (\textit{Data Ecologies}, 2024) and Vivien Jiaqian Zhu’s dramaturgical critique of the \textbf{corporate stage} (\textit{The Evolving Landscape of Workplace}, 2027). 

We argue that contemporary algorithmic governance operates through a closed cybernetic feedback loop across two distinct registers:
\begin{enumerate}[label=\arabic*.]
    \item \textbf{The Outer Architectural Circuit (Aho):} Code-level parameters and dynamic telemetry controls that restrict user choices at the interface layer, systematically stripping software of its \textit{interpretive flexibility}.
    \item \textbf{The Inner Subjective Circuit (Zhu):} Dramaturgical dynamics that force workers to perform productivity on digital stages while anxiously monitoring their own live metric proxies on automated dashboards.
\end{enumerate}

% --- Section 2 ---
\section{The Outer Circuit: Architectural Enclosure and Telemetry Control}

Aho’s framework treats software platforms not as neutral productivity tools, but as active digital ecologies governed through boundary control and telemetry ingestion. In the context of the quantified workplace, enterprise software functions as a sensor-saturated habitat designed to capture every ambient trace of human activity---from keystroke volume and mouse trajectory to active window latency and message dispatch frequency.

\subsection{Eradicating Interpretive Flexibility}
In Science and Technology Studies (STS), \textit{interpretive flexibility} refers to the capacity of human users to re-interpret, adapt, or manipulate technology in ways unanticipated by its original designers. In the modern workplace, workers frequently exercise interpretive flexibility to reclaim autonomy, protect themselves from exhaustion, or obscure periods of rest. These adaptations include:
\begin{itemize}
    \item \textbf{Hardware Interventions:} Utilizing physical mouse jigglers or mechanical weight devices to simulate continuous activity.
    \item \textbf{Software Workarounds:} Scripting automated status updates, running background keep-alive macros, or employing split-screen configurations to obscure idle time.
    \item \textbf{Strategic Communication:} Scheduling message dispatches outside working hours to fabricate an aura of continuous availability.
\end{itemize}

\subsection{Software Shepherding in Action}
Under Aho’s model, platform developers use background telemetry to identify these unprescribed adaptations in real time. Rather than issuing manual disciplinary warnings, system designers execute \textbf{software shepherding}: deploying automated code updates, interface speed bumps, and administrative patches that close technical loopholes. 

\begin{quote}
\textit{Software shepherding operates by altering the technical boundaries of the interface itself---eliminating the physical or computational possibility of alternative user behavior.}
\end{quote}

When an enterprise update automatically detects hardware-simulated mouse movements and adjusts status indicators to ``Away,'' or when an interface automatically logs out a user during a period of window inactivity, the platform exercises architectural enclosure. The software shepherd forces the worker back into prescribed operational pathways without engaging in interpersonal negotiation.

% --- Section 3 ---
\section{The Inner Circuit: The Corporate Stage and the Paradox of Spectatorship}

While Aho identifies how code enforces boundaries from the outside, Zhu accounts for how workers experience and navigate these boundaries from within. Adapting Erving Goffman’s dramaturgical theory to digital environments, Zhu conceptualizes the hybrid workspace as a \textbf{corporate stage}.

\subsection{The Collapse of the Backstage}
In traditional physical offices, clear boundaries separate the ``frontstage'' (where formal professional performance occurs) from the ``backstage'' (break rooms, hallway conversations, private desks where workers relax and drop their professional personas). Hybrid and remote platforms collapse this distinction. Because physical presence is replaced by digital telemetry, the worker’s home workspace is subsumed into a continuous, stylized frontstage performance. To be offline, idle, or unmeasured is to cease existing within the organizational structure.

\subsection{The Paradox of Spectatorship}
This collapsed environment gives rise to what Zhu terms the \textbf{paradox of spectatorship}. Workers in remote environments exist in a state of dual visibility: they are hyper-visible to automated tracking infrastructure yet interpersonally invisible to human managers and peers. To bridge this gap, workers must constantly curate their digital presence, transforming routine tasks into stylized performances of responsiveness and activity.

Simultaneously, workers become passive, anxious spectators of their own quantified data proxies. As enterprise platforms expose productivity dashboards, response-time scores, and activity percentages directly to the employee, the worker is forced to look into a ``digital mirror.'' They do not merely perform labor; they continuously monitor the computational shadow of their labor to ensure it aligns with algorithmic expectations.

% --- Section 4 ---
\section{Synthesis: The Co-Production of Automated Self-Discipline}

Synthesizing Aho’s structural model with Zhu’s dramaturgical framework reveals that algorithmic management operates as an integrated system of \textbf{co-produced automated self-discipline}. The outer architectural circuit and the inner subjective circuit are not parallel processes; they are mutually reinforcing halves of a single cybernetic mechanism.

\begin{table}[h]
\centering
\caption{Synthesis of the Dual-Circuit Governance Model}
\label{tab:dual_circuit}
\small
\begin{tabular}{p{3cm} p{6cm} p{6cm}}
\toprule
\textbf{Dimension} & \textbf{Outer Architectural Circuit (Aho)} & \textbf{Inner Subjective Circuit (Zhu)} \\
\midrule
\textbf{Primary Domain} & Technical Code \& Interface Architecture & Subjective Experience \& Dramaturgical Performance \\
\textbf{Key Mechanism} & Telemetry Ingestion \& Software Shepherding & Dashboard Monitoring \& Spectator Curation \\
\textbf{Target of Control} & Interpretive Flexibility \& User Workarounds & Worker Agency \& Performative Status \\
\textbf{Operational Mode} & Structural Barrier (The Cage) & Self-Monitoring Mirror (The Proxy) \\
\textbf{Behavioral Result} & Code-Enforced Limitations & Preemptive Self-Correction \\
\bottomrule
\end{tabular}
\end{table}

\subsection{Preemptive Self-Shepherding}
The key link between these circuits is \textbf{preemptive self-shepherding}. Because workers are positioned as anxious spectators of their own metric proxies, they do not wait for an automated soft-lockout, an administrative flag, or a system patch to adjust their behavior. The constant psychological friction of dashboard self-monitoring forces the worker to internalize the platform's rules.

Workers modify their physical inputs, work rhythms, and communication styles in advance to ensure their quantified proxy remains within acceptable parameters:
\begin{itemize}
    \item Typing periodically during document review to prevent the idle-counter from triggering.
    \item Keeping communication channels open during breaks to avoid ``Away'' status transitions.
    \item Altering daily task sequences to prioritize tasks that produce clear telemetry traces over qualitative tasks that do not.
\end{itemize}

Aho's outer circuit enforces what the interface permits, while Zhu's inner circuit ensures the worker actively maintains the lock by policing their own body. Algorithmic management achieves total enclosure by conflating labor performance with metric stewardship.

% --- Section 5 ---
\section{Methodological and Empirical Framework}

Investigating this dual-circuit architecture requires a bimodal empirical strategy capable of capturing both technical code modifications and subjective worker experiences.

\subsection{Technical \& Interface Auditing (Outer Circuit)}
To map software shepherding, scholars must employ critical data auditing techniques:
\begin{itemize}
    \item \textbf{Telemetry Packet Logging:} Inspecting API endpoints and network traffic generated by enterprise software to determine what ambient data is captured, at what frequency, and under what conditions.
    \item \textbf{Patch \& Version Analysis:} Documenting update histories across platforms (e.g., Teams, Teramind, Hubstaff) to trace how developers systematically identify and eliminate user workarounds.
    \item \textbf{Dark Pattern Identification:} Cataloging interface design choices that introduce friction into privacy-preserving or autonomy-seeking user configurations.
\end{itemize}

\subsection{Digital Ethnography \& Phenomenological Inquiries (Inner Circuit)}
To examine the paradox of spectatorship, empirical research must capture the lived reality of remote workers:
\begin{itemize}
    \item \textbf{Metric Diary Studies:} Having workers log their daily interactions with productivity dashboards, recording moments of anxiety, performative status curation, and physical self-correction.
    \item \textbf{Qualitative In-Depth Interviews:} Utilizing semi-structured interview protocols focused on hyper-visibility, backstage collapse, and the emotional labor of maintaining digital active indicators.
    \item \textbf{Triangulated Observation:} Comparing worker self-reports with interface telemetry logs to isolate instances of preemptive self-shepherding.
\end{itemize}

\subsection{Future Research: Generative AI and Cognitive Telemetry}
As enterprise platforms integrate generative AI assistants (e.g., Microsoft 365 Copilot, Google Workspace AI), the dual circuit of workplace governance is undergoing a qualitative shift: from \textit{activity-based telemetry} (keystrokes, mouse movement, idle alerts) to \textit{cognitive and affective telemetry} (prompt histories, document synthesis speeds, communication tone, and real-time meeting transcripts). This expands Aho's software shepherding from regulating physical interaction inputs to steering thought processes and writing styles, while deepening Zhu's paradox of spectatorship through automated evaluations of worker sentiment and cognitive efficiency.

\clearpage

% --- Appendices ---
\appendix

\section{Inter-Coder Reliability (ICR) Protocol and Statistical Validation}

This appendix outlines the qualitative coding workflow, statistical agreement metrics, and discrepancy reconciliation procedures implemented to establish coding reliability and codebook stability across the interview corpus.

\subsection{Two-Coder Calibration and Blinded Coding Workflow}
Qualitative analysis was executed by two independent researchers: the Primary Investigator (Coder A) and an independent qualitative specialist (Coder B). The ICR process followed a three-tier operational pipeline:
\begin{enumerate}
    \item \textbf{Phase 1 (Codebook Orientation \& Pilot Calibration):} Both researchers thoroughly reviewed the structured codebook, inclusion/exclusion guidelines, and exemplar transcript excerpts. Coders completed a pilot trial on a 10\% transcript subset, followed by a joint calibration meeting to standardize code boundary definitions.
    \item \textbf{Phase 2 (Blinded Independent Coding):} A randomized 20\% sample of transcript segments ($N = 100$ discrete coded units) was selected for formal reliability auditing. Both coders independently applied tags to discrete text units using CAQDAS software (NVivo/MAXQDA) while remaining blinded to each other's tagging outputs.
    \item \textbf{Phase 3 (Matrix Merging \& Audit Trail):} Independent coding spreadsheets were compiled into a master comparative matrix to isolate exact tag matches, partial overlaps, and structural disagreements.
\end{enumerate}

\subsection{Statistical Calculation of Inter-Coder Agreement}
To account for agreement occurring by random chance across a 10-item sub-code taxonomy, inter-coder reliability was evaluated using Cohen’s Kappa ($\kappa$):

\begin{equation}
\kappa = \frac{p_o - p_e}{1 - p_e}
\end{equation}

Where $p_o$ represents the observed proportion of agreement and $p_e$ represents the expected proportion of chance agreement. Based on the evaluation of $N = 100$ discrete coded units across the 10 sub-codes:
\begin{itemize}
    \item Observed Proportional Agreement ($p_o$): $0.80$ ($80.0\%$)
    \item Expected Chance Agreement ($p_e$): $0.10$
    \item Calculated Coefficient:
\end{itemize}

\begin{equation}
\kappa = \frac{0.80 - 0.10}{1 - 0.10} = \frac{0.70}{0.90} \approx 0.778
\end{equation}

In accordance with Landis and Koch’s (1977) evaluation benchmarks, $\kappa = 0.778$ confirms \textbf{substantial inter-coder agreement} ($0.61 \le \kappa \le 0.80$), validating the stability and reproducibility of the coding schema.

\subsection{Discrepancy Resolution Protocol}
All coding discrepancies ($20\%$ of evaluated segments) were resolved through a standardized reconciliation process:
\begin{itemize}
    \item \textbf{Primary Intent Rule:} In edge cases where segments exhibited overlap between technical infrastructure awareness (\texttt{TELE-SENS}) and performative visibility (\texttt{STAGE-VISI} / \texttt{STAGE-PRES}), coders evaluated the participant’s primary communicative intent in the surrounding narrative block.
    \item \textbf{Consensus Reconciliation:} Segments with conflicting tags were re-analyzed against the explicit exclusion criteria in the codebook during post-coding debriefs until $100\%$ final consensus was reached.
    \item \textbf{Audit Trail Logging:} All decision rules and boundary refinements established during consensus sessions were recorded in a qualitative audit trail to preserve analytical transparency.
\end{itemize}

\clearpage

\section{Statement of Positionality and Reflexivity}

As qualitative researchers investigating algorithmic management and workplace surveillance, our analytical lens is inherently shaped by our institutional locations, theoretical orientations, and lived experiences within modern knowledge-work ecosystems. As academic researchers operating in hybrid university environments, we routinely interact with enterprise collaboration software, automated presence indicators, and background telemetry. Recognizing that our intimate familiarity with these digital environments could introduce unexamined assumptions regarding worker fatigue, autonomy, and technological enclosure, we executed systematic reflexive protocols throughout data collection and analysis to safeguard analytical trustworthiness.

Our research team synthesizes theoretical frameworks from Science and Technology Studies (STS) and dramaturgical organizational sociology. While this orientation provided the conceptual architecture for examining software shepherding and spectator performance, we remained vigilant against deductively forcing participant narratives into pre-conceived categories. To mitigate researcher confirmation bias, we maintained three operational safeguards:
\begin{itemize}
    \item \textbf{Pre-Analytical Epoché:} Prior to coding, researchers authored reflexive positionality memos documenting personal biases and assumptions regarding remote work oversight.
    \item \textbf{Peer Debriefing Audits:} Bi-weekly calibration sessions were conducted to interrogate coding splits, ensuring sub-codes grounded themselves inductively in transcript excerpts rather than theoretical expectations.
    \item \textbf{Negative Case Integration:} We deliberately sought out, coded, and analyzed participant accounts that contradicted theories of algorithmic control, such as expressions of worker indifference or unmonitored autonomy.
\end{itemize}

By documenting our positionality and maintaining these analytical checks, we ensure that our findings authentically reflect the complex, heterogeneous realities of contemporary remote labor.

\clearpage

\section{REFI-QDA Standard XML Codebook (\texttt{.qdc})}

The XML structure below provides the complete, machine-readable codebook taxonomy formatted under the REFI-QDA standard for direct import into QDA software (NVivo, MAXQDA, ATLAS.ti, Dedoose).

\begin{lstlisting}
<?xml version="1.0" encoding="utf-8"?>
<CodeBook xmlns="urn:QDA-XML:codebook:1.0" origin="Aho-Zhu Framework" creationDateTime="2026-08-07T08:45:00Z">
  <Codes>
    <!-- Parent Code 1 -->
    <Code guid="c1000000-0000-0000-0000-000000000001" name="1. Ambient Telemetry [TELE]" isConceptual="true">
      <Description>Infrastructural capture of background user data and habitat saturation.</Description>
      <SubCodes>
        <Code guid="c1100000-0000-0000-0000-000000000001" name="1.1 Sensor Awareness [TELE-SENS]">
          <Description>INCLUSION: Recognition of background tracking (idle counters, keystrokes, webcam flags). EXCLUSION: Active monitoring of self-metrics (SPEC-DASH) or tricking sensors (SHEP-FLEX). EXEMPLAR: "I'm always conscious that Teams flips my status icon..."</Description>
        </Code>
        <Code guid="c1200000-0000-0000-0000-000000000001" name="1.2 Habitat Saturation [TELE-HABI]">
          <Description>INCLUSION: Remote space feeling integrated into corporate data infrastructure. EXCLUSION: General emotional burnout lacking reference to software (STAGE-EROD). EXEMPLAR: "My home office stopped feeling like my personal house..."</Description>
        </Code>
      </SubCodes>
    </Code>

    <!-- Parent Code 2 -->
    <Code guid="c2000000-0000-0000-0000-000000000002" name="2. Interface Shepherding [SHEP]" isConceptual="true">
      <Description>Code-level boundaries, user workarounds, and algorithmic counter-measures.</Description>
      <SubCodes>
        <Code guid="c2100000-0000-0000-0000-000000000002" name="2.1 Interpretive Flexibility [SHEP-FLEX]">
          <Description>INCLUSION: Hacks, physical tools (mouse jigglers), or scripts to obscure idle time. EXCLUSION: Standard task prioritization or sanctioned productivity shortcuts. EXEMPLAR: "I bought a mechanical mouse jiggler on Amazon..."</Description>
        </Code>
        <Code guid="c2200000-0000-0000-0000-000000000002" name="2.2 Algorithmic Counter-Measures [SHEP-PATC]">
          <Description>INCLUSION: Software updates, patches, or system locks disabling user workarounds. EXCLUSION: General IT software bugs or hardware failures. EXEMPLAR: "After the latest security patch, my keep-alive script stopped working..."</Description>
        </Code>
      </SubCodes>
    </Code>

    <!-- Parent Code 3 -->
    <Code guid="c3000000-0000-0000-0000-000000000003" name="3. The Corporate Stage [STAGE]" isConceptual="true">
      <Description>Dramaturgical dynamics of remote work, visibility, and backstage collapse.</Description>
      <SubCodes>
        <Code guid="c3100000-0000-0000-0000-000000000003" name="3.1 Performative Presence [STAGE-PRES]">
          <Description>INCLUSION: Stylized behaviors executed solely to signal availability to peers/managers. EXCLUSION: Routine, task-focused professional communication. EXEMPLAR: "I deliberately hold off on sending non-urgent Slack messages until 6:00 PM..."</Description>
        </Code>
        <Code guid="c3200000-0000-0000-0000-000000000003" name="3.2 Backstage Collapse [STAGE-EROD]">
          <Description>INCLUSION: Personal home spaces or breaks invaded by digital visibility requirements. EXCLUSION: Complaints about physical office environments or commutes. EXEMPLAR: "There is no 'off' button anymore..."</Description>
        </Code>
        <Code guid="c3300000-0000-0000-0000-000000000003" name="3.3 Paradoxical Visibility [STAGE-VISI]">
          <Description>INCLUSION: Being hyper-observed by software logging while feeling interpersonally isolated. EXCLUSION: General workplace loneliness lacking connection to software metrics. EXEMPLAR: "The system tracks every second... but my manager hasn't spoken to me in weeks."</Description>
        </Code>
      </SubCodes>
    </Code>

    <!-- Parent Code 4 -->
    <Code guid="c4000000-0000-0000-0000-000000000004" name="4. Spectator Performance [SPEC]" isConceptual="true">
      <Description>Self-monitoring, preemptive compliance, and metric displacement.</Description>
      <SubCodes>
        <Code guid="c4100000-0000-0000-0000-000000000004" name="4.1 Dashboard Self-Monitoring [SPEC-DASH]">
          <Description>INCLUSION: Checking or stress associated with viewing live productivity percentages/proxies. EXCLUSION: Checking personal work calendars or standard project task boards. EXEMPLAR: "Every afternoon at 4:00 PM, I open my Teramind dashboard..."</Description>
        </Code>
        <Code guid="c4200000-0000-0000-0000-000000000004" name="4.2 Preemptive Self-Shepherding [SPEC-PREE]">
          <Description>INCLUSION: Modifying physical inputs in advance to avoid triggering an idle alert. EXCLUSION: Post-hoc corrections made after receiving a human reprimand. EXEMPLAR: "I force myself to tap the spacebar or scroll randomly every two minutes..."</Description>
        </Code>
        <Code guid="c4300000-0000-0000-0000-000000000004" name="4.3 Metric Displacement [SPEC-DISP]">
          <Description>INCLUSION: Prioritizing low-value tasks that generate telemetry over high-value qualitative work. EXCLUSION: Standard procrastination or personal distractions. EXEMPLAR: "I spent two hours answering quick, trivial emails just to keep my activity bar green..."</Description>
        </Code>
      </SubCodes>
    </Code>
  </Codes>
</CodeBook>
\end{lstlisting}
\bibliographystyle{plain}
\bibliography{references}
\end{document}